% ============================================================
\section{Targets and Certifications}
\label{sec:targets}
% ============================================================

\begin{quote}\small
\textbf{Goal:} Catalog the structural and objective targets in Coherence Calculus.\\
\textbf{Structural targets} are exact certification statements: properties that can be certified directly from the active interfaces, either on the whole visible state or on a selected workload.\\
\textbf{Objective targets} are optimization problems: find a projection that minimizes/maximizes a functional.
\end{quote}

% ------------------------------------------------------------
\subsection{Structural Targets (Iff Theorems)}
\label{subsec:structural-targets}
% ------------------------------------------------------------

Structural targets are exact certifications on the active interface surface.
Some are equivalences, while others are one-way sufficient criteria.
No approximation or optimization is required; only verification.

\paragraph{A) Perfect Coherence (Zero Residual)}

\begin{definition}[Perfect coherence]
\label{def:perfect-coherence}\lean{CoherenceCalculus.hasPerfectCoherence}
An operator $A$ has \emph{perfect coherence} with respect to projection $P$ if the Schur correction vanishes:
\[
\CCSigmaP{P}{z} = 0 \quad \text{for all $z$ where defined}.
\]
Equivalently, $\CCAeff{A}{P}(z) = \CCBlockPP{A}$ for all $z$.
\end{definition}

\begin{theorem}[Cross-block vanishing implies perfect coherence]
\label{thm:block-diagonal-sufficient}\lean{CoherenceCalculus.zeroPQ_implies_perfect, CoherenceCalculus.zeroQP_implies_perfect}
Let $A$ be a bounded operator on $\cH$ and $P$ an orthogonal projection with $\CCComp = I - P$.
If either cross-block vanishes, then the Schur correction vanishes identically:
\[
\CCBlockPQ{A} = 0 \text{ or } \CCBlockQP{A} = 0
\quad \Longrightarrow \quad
\CCSigmaP{P}{z} = 0 \text{ for all $z$ where defined}.
\]
\end{theorem}

\begin{proof}
If $\CCBlockPQ{A} = 0$ or $\CCBlockQP{A} = 0$, then
\[
\CCSigmaP{P}{z}
=
\CCBlockPQ{A} (I - z\CCBlockQQ{A})^{-1} \CCBlockQP{A}
=
0.
\]
\end{proof}

\begin{remark}[Exact coefficient criterion for perfect coherence]
\label{rem:perfect-coherence-coefficients}
\lean{schur_transport_jet, schur_transport_jet_order}
The converse of Theorem~\ref{thm:block-diagonal-sufficient} is false in general:
perfect coherence does \emph{not} force one of the two cross-blocks to vanish.
If \(\CCSigmaP{P}{z}=0\) for all \(z\) in a neighborhood of \(0\), then the
resolvent expansion gives the coefficient-vanishing condition
\[
\CCBlockPQ{A} (\CCBlockQQ{A})^k \CCBlockQP{A}=0
\qquad\text{for every }k\ge 0.
\]
Thus cross-block vanishing is a sufficient certificate for perfect coherence,
but not a necessary one.
\end{remark}

\begin{corollary}[Commutant characterization]
\label{cor:commutant-perfect}\lean{CoherenceCalculus.commutant_implies_perfect}
If $[A, P] = 0$ (equivalently, $\CCBlockPQ{A} = \CCBlockQP{A} = 0$), then $A$ has perfect coherence with respect to $P$.
\end{corollary}

\begin{quote}\small
\textbf{Certification:} A fast sufficient certificate for perfect coherence is
to verify $\CCBlockPQ{A} = 0$ or $\CCBlockQP{A} = 0$.
If one cross-block vanishes, no Schur correction is needed.
The exact necessary-and-sufficient criterion is the coefficient vanishing
condition of Remark~\ref{rem:perfect-coherence-coefficients}.
\end{quote}

\paragraph{B) Task-Perfect Coherence for a Selected Workload}

Section~\ref{sec:observable-relative-closure} introduces the workload-relative
return transfer
\[
\widetilde{\mathfrak R}_{\Omega}(P,A;z)
\]
for an observable package \(\Omega=(S,V,O)\). This fixed-source version keeps
the workload domain equal to \(S\) while
Proposition~\ref{prop:fixed-source-return-restricts} shows that it restricts
to the pointwise visible-workload return on \(S\cap P\cH\). By
Definition~\ref{def:task-perfect-coherence}, the pair \((A,P)\) is
\emph{task-perfect} for \(\Omega\) exactly when this transfer vanishes.
Corollary~\ref{cor:perfect-implies-task-perfect} shows that whole-state perfect
coherence implies task-perfect coherence, but not conversely: the workload may
ignore returned directions that remain present in the full state.
Section~\ref{sec:observable-relative-closure} also packages the source-visible
loss, direct observable transfer, and fixed-source Schur return into the
task-visible defect package \(\mathcal E_{\Omega}(P,A;z)\), so one can
distinguish source-faithful, one-step faithful, task-perfect, and task-exact
cuts without conflating them.

\begin{quote}\small
\textbf{Certification:} Check that
\(\widetilde{\mathfrak R}_{\Omega}(P,A;z)=0\) for all \(z\) where defined.
Equivalently, after projecting the workload source into the visible cut, the
selected workload sees no returned residual forcing.
\end{quote}

\paragraph{C) Solver Exactness}

\begin{definition}[Solver exactness]
\label{def:solver-exactness}\lean{CoherenceCalculus.isSolverExact}
A Hilbert complex $(H^k, D_k)$ is \emph{solver-exact at degree $k$} if
\[
\CCCoh{k}(D) = 0 \quad \Longleftrightarrow \quad \Ker(D_k) = \Img(D_{k-1}).
\]
\end{definition}

\begin{quote}\small
\textbf{Certification:} Check $\CCCoh{k}(D) = 0$.
If satisfied, the Green operator $\CCGreen{k} = \CCLap{k}^{-1}$ exists and provides canonical integration via $\CCHom{k} = D_{k-1}^* \CCGreen{k}$.
See Section~\ref{subsec:solver-layer} for details.
\end{quote}

\paragraph{D) Realization Exactness and Goal Certificates}

Section~\ref{sec:realization-calculus} adds the realization-side certificate
surface. A realization chain carries the exact residual ledger
\(\mathfrak C_h\) of
Definition~\ref{def:typed-crime-ledger}. By
Theorem~\ref{thm:exact-residual-ledger}, the difference between the fully
realized residual and the pulled-back continuum residual is exactly the sum of
those typed crimes. Definition~\ref{def:realization-exactness} then names the
strict structural target: the realization is exact iff every declared crime
vanishes. The same section also provides the dual-weighted
goal-error certificate of Theorem~\ref{thm:exact-goal-error-ledger-identity},
which decomposes output error against the same ledger.

\begin{quote}\small
\textbf{Certification:} Check that each declared crime in \(\mathfrak C_h\)
vanishes. Then the host realizes the pulled-back continuum residual exactly.
If only a selected output matters, solve the adjoint goal problem and evaluate
the dual-weighted ledger to see which crimes remain visible to that output.
\end{quote}

\paragraph{E) Holographic Saturation}

\begin{definition}[Holographic saturation (repeated)]
\label{def:holographic-saturation-targets}\lean{CoherenceCalculus.isHolographicallySaturated}
A complex with boundary map $R_k$ is \emph{holographically saturated at degree $k$} if
\[
\CCCohRel{k}{D,R} = 0 \quad \Longleftrightarrow \quad \Ker(R_k) = \Img(D_{k-1}).
\]
\end{definition}

\begin{quote}\small
\textbf{Certification:} Check that $\Ker(R_k) = \Img(D_{k-1})$.
If satisfied, zero boundary data implies the bulk state is pure gauge (exact).
See Definition~\ref{def:holographic-saturation}.
\end{quote}

% ------------------------------------------------------------
\subsection{Objective Targets (Optimization Interface)}
\label{subsec:objective-targets}
% ------------------------------------------------------------

Objective targets involve finding a projection $P$ that optimizes some functional.
Layer-1 provides the existence theorems and stationarity conditions; Layer-2 (applications) chooses the specific objective.

\paragraph{The Grassmannian of Projections}

\begin{definition}[Grassmannian]
\label{def:grassmannian}\lean{CoherenceCalculus.isInGrassmannian, CoherenceCalculus.GrassmannianElem}
For a finite-dimensional Hilbert space $\cH$, the \emph{Grassmannian of rank-$r$ projections} is
\[
\CCGrass{r} := \{ P \in \cB(\cH) : P^2 = P = P^*, \; \rank(P) = r \}.
\]
This is a compact smooth manifold of dimension $r(n - r)$, where $n = \dim \cH$.
\end{definition}

\paragraph{Generic Objective Functional}

\begin{definition}[Design objective]
\label{def:design-objective}\lean{CoherenceCalculus.DesignObjective}
A \emph{design objective} for the inverse split problem is a continuous functional
\[
\CCDesign{P} := \|\CCSigmaP{P}{z} - \Sigma^\star\|^2 + \lambda \, \mathrm{Reg}(P),
\]
where:
\begin{itemize}[itemsep=0.2em]
\item $\Sigma^\star$ is a target Schur correction (possibly zero for perfect coherence).
\item The norm is Frobenius/Hilbert--Schmidt in finite dimensions.
\item $\mathrm{Reg}(P)$ is a regularizer (e.g., $\|[A, P]\|^2$, $\Tr(P)$, or $0$).
\item $\lambda \ge 0$ is a regularization weight.
\end{itemize}
\end{definition}

\begin{definition}[Workload-relative design objective]
\label{def:workload-design-objective}\lean{CoherenceCalculus.WorkloadDesignObjective}
Let \(\Omega=(S,V,O)\) be a finite-dimensional observable package, let
\(m\in\bbN\), and let \(\mu,\lambda\ge 0\). A \emph{workload-relative design
objective} is
\[
J_{\Omega,m}(P)
:=
\|O(I-P)\big|_{S}\|^2
+
\mu\, e_m\!\bigl(\widetilde{\mathfrak R}_{\Omega}(P,A;z)\bigr)^2
+
\lambda\,\mathrm{Reg}(P),
\]
where \(e_m\) is the best \(m\)-dimensional correction error from
Definition~\ref{def:best-m-correction-error}.
\end{definition}

\paragraph{Interpretation.}
The first term measures what the cut loses directly on the selected workload.
The second measures how much of the returned residual remains after allowing at
most \(m\) correction coordinates. By
Corollary~\ref{cor:optimal-returned-correction-for-workload}, in a
finite-dimensional Hilbert package this second term is exactly the square of
the \((m+1)\)th singular value of the fixed-source observable return transfer.
Equivalently, by Theorem~\ref{thm:observable-closure-invariant-package}, it is
the \((m+1)\)th eigenvalue of the observable closure Gramian
\(G_{\Omega}(P,A;z)\). The third is any geometric or structural regularizer
already present in the operator-design problem.

\paragraph{Shadow-transport objectives.}
\label{ex:shadow-transport-objectives}
For a projection $P$, let
\[
B_P:=(I-P)\,d\,P,
\qquad
T_P:=B_P^*B_P.
\]
Natural objective functionals built from the shadow-transport invariant family
include:
\[
\begin{aligned}
J_{\mathrm{tot}}(P) &:= \Tr(T_P), \\
J_k(P) &:= \Tr(T_P^k) \qquad (k\ge 2), \\
J_{\det,t}(P) &:= \det(I+tT_P), \\
J_{\max}(P) &:= \|T_P\|.
\end{aligned}
\]
These measure, respectively, total leakage, higher spectral moments, the
determinant polynomial at fixed parameter $t$, and maximal leakage.

\paragraph{Existence of Minimizers}

\begin{theorem}[Existence of minimizers]
\label{thm:existence-minimizers}\lean{CoherenceCalculus.existence_of_minimizers}
Let $\cH$ be finite-dimensional, $A$ a bounded operator, and $\CCDesign{P}$ a continuous design objective with invertibility margin maintained on $\CCGrass{r}$.
Then:
\begin{enumerate}[label=(\roman*)]
\item $\CCGrass{r}$ is compact.
\item $\CCDesign{\cdot}$ is continuous on $\CCGrass{r}$.
\item A minimizer $P^\star \in \CCGrass{r}$ exists.
\end{enumerate}
\end{theorem}

\begin{proof}
(i) $\CCGrass{r}$ is a closed subset of the unit ball in $\cB(\cH)$ (since $\|P\| \le 1$ for projections), hence compact in finite dimensions.

(ii) Continuity of $P \mapsto \CCSigmaP{P}{z}$ follows from continuity of the block extraction $P \mapsto PAP$, etc., and continuity of the inverse $(I - z\CCBlockQQ{A})^{-1}$ under the invertibility margin assumption.

(iii) Continuous function on a compact set attains its minimum.
\end{proof}

\begin{corollary}[Certification without construction]
\label{cor:certification-existence}\lean{CoherenceCalculus.certification_without_construction}
Under the hypotheses of Theorem~\ref{thm:existence-minimizers}, an optimal projection $P^\star$ exists even if we cannot construct it explicitly.
This provides a ``certification without construction'' guarantee.
\end{corollary}

\begin{corollary}[Workload-relative minimizers]
\label{cor:workload-relative-minimizers}\lean{CoherenceCalculus.workload_relative_minimizers}
Let \(\cH\) be finite-dimensional, let \(A\) be bounded, and let
\(\Omega=(S,V,O)\) be a finite-dimensional observable package. Assume the
invertibility margin for the Schur correction is maintained on \(\CCGrass{r}\).
Then every workload-relative design objective \(J_{\Omega,m}\) from
Definition~\ref{def:workload-design-objective} attains a minimum on
\(\CCGrass{r}\).
\end{corollary}

\begin{proof}
The map \(P\mapsto O(I-P)\big|_S\) is continuous. Under the maintained
invertibility margin, Theorem~\ref{thm:existence-minimizers} gives continuity
of \(P\mapsto \CCSigmaP{P}{z}\). Proposition~\ref{prop:fixed-source-return-continuity}
therefore yields continuity of
\(P\mapsto \widetilde{\mathfrak R}_{\Omega}(P,A;z)\). Proposition~\ref{prop:best-m-correction-distance}
shows that \(e_m\) is Lipschitz, so
\[
P\longmapsto e_m\!\bigl(\widetilde{\mathfrak R}_{\Omega}(P,A;z)\bigr)
\]
is continuous as well. Therefore \(J_{\Omega,m}\) is continuous on the compact
Grassmannian \(\CCGrass{r}\), and the minimum exists by
Theorem~\ref{thm:existence-minimizers}.
\end{proof}

\paragraph{Stationarity Condition}

The stationarity condition for minimizing $\CCDesign{P}$ on $\CCGrass{r}$ involves the tangent space of the Grassmannian.
See Section~\ref{sec:projector-calculus} for the projector differential calculus and the derivation of:
\[
P^\star \text{ is stationary} \quad \Longleftrightarrow \quad [\nabla \CCDesign{P^\star}, P^\star] = 0,
\]
where $\nabla \CCDesign{P}$ is the Euclidean gradient of the objective.

% ------------------------------------------------------------
\subsection{Summary Table}
\label{subsec:target-summary}
% ------------------------------------------------------------

\begin{center}
\renewcommand{\arraystretch}{1.3}
\begin{tabular}{@{}lll@{}}
\textbf{Target} & \textbf{Type} & \textbf{Certification} \\
\hline
Perfect coherence & Structural & $\CCBlockPQ{A} = 0$ or $\CCBlockQP{A} = 0$ \\
Task-exact cut & Structural & source-faithful + one-step faithful + task-perfect \\
Task-perfect coherence & Structural & $\widetilde{\mathfrak R}_{\Omega}(P,A;z)=0$ \\
Solver exactness & Structural & $\CCCoh{k}(D) = 0$ \\
Holographic saturation & Structural & $\Ker(R_k) = \Img(D_{k-1})$ \\
\hline
Optimal projection & Objective & $[\nabla J, P] = 0$ (stationarity) \\
Optimal workload cut & Objective & minimize $J_{\Omega,m}(P)$ \\
Existence guarantee & Objective & Compactness of $\CCGrass{r}$ \\
\end{tabular}
\end{center}
